\section{A stationary current and its line trajectories}
\label{stat:section}

We retain the restricted HNN splitting and notation of
Theorem~\ref{ann:combination}.  Thus
\[
 G=\langle H,t\mid tAt^{-1}=B\rangle,\qquad P=A,
 \qquad I=A\cap B,
\]
and $i,j:I\longrightarrow P$ are inclusion and conjugation by $t^{-1}$,
with respective images $I,J$.  In this section $G$ is assumed to contain
no Baumslag--Solitar subgroup.  Our first goal is to replace the bounded
sequence of currents supplied by Proposition~\ref{ann:trajectory} with
one nonzero current satisfying
\[
                    i_*\nu=j_*\nu=\mu.
\]
We will then interpret this equality as a measure-preserving transfer
of almost every line along a unique path in the Bass--Serre tree
$\mathcal T$.

\subsection{Periodic axes and a finite set of exceptional roots}

A \emph{root of an axis} in the Cayley tree of $P$ means a generator of
its stabilizer in $P$.  Such a generator is not a proper power; it need
not belong to a free basis.  We consider roots up to conjugacy and
inversion, or equivalently consider the corresponding conjugacy
classes of maximal cyclic subgroups.

\begin{lemma}\label{stat:atomic-axis}
Every atom of a free-group current is a periodic axis.  If the mass of
one oriented axis is $b>0$ and its root is $r$, then its translate orbit,
together with the reversed orientations, contributes
$b|r|_{\mathrm{cyc}}$ to the current norm.
\end{lemma}
\begin{proof}
Translate successive vertices of an atomic line to the identity.  All
resulting lines have the same positive mass and pass through the
identity.  That set of lines is compact and has finite measure, so two
of the translates coincide.  A nontrivial element of the free group
therefore preserves the original line.  An element reversing a line
would fix a vertex or invert an edge, which is impossible for the free
action on the Cayley tree.  The line stabilizer consequently acts by
translations and is infinite cyclic.

For its root $r$, the number of distinct translates of one orientation
through the identity is $|r|_{\mathrm{cyc}}$.  This is the number of
vertices in the quotient axis by $\langle r\rangle$.  Flip invariance
gives the same mass to the other orientation.  The definition of the
norm as half the mass of the lines through the identity gives the
claimed contribution.
\end{proof}

Recall that a transition type $\tau=(\epsilon,\delta,d)$ consists of
an intersection group $K_\tau$ and two injections
\[
 \theta_{\tau,0},\theta_{\tau,1}:K_\tau\longrightarrow P.
\]
There are finitely many types in Proposition~\ref{ann:trajectory}.
Their traversal signs and double-coset representatives are part of the
type.  A noncyclic type must be aligned: its two signs agree, its
representative is $1$, and its two maps are $i,j$, in one of the two
orders.

\begin{lemma}\label{stat:finite-roots}
There is a finite set $\mathcal D$ of root classes with the following
properties.
\begin{enumerate}[label=\textup{(\roman*)}]
\item The root of either image of a nontrivial element of a cyclic
transition group belongs to $\mathcal D$.
\item If a periodic axis is carried by $p\Lambda E$ and $q\Lambda E$
for distinct cosets $pE\ne qE$, where $E=I$ or $E=J$, its root
belongs to $\mathcal D$.
\item If an axis root in a noncyclic transition group maps to a proper
power under either transition injection, the root of that image
belongs to $\mathcal D$.
\end{enumerate}
Here ``carried'' means that both endpoints lie in the indicated
translate of the limit set.
\end{lemma}
\begin{proof}
For each nontrivial cyclic transition group, include the roots of the
images of one generator under both injections.  This makes a finite
set and proves~(i).

For $E=I,J$, the strong inertia inequality with test subgroup $E$
implies that every off-diagonal intersection
$E\cap xEx^{-1}$, $x\notin E$, is at most cyclic: the diagonal term
already contributes $\bar r(E)$.  By
Lemma~\ref{pre:double-coset-finiteness}, only finitely many double
cosets $ExE$ have nontrivial intersection.  For each such off-diagonal
double coset, include the root in $P$ of a generator of its
intersection.

Suppose that an axis is carried by $p\Lambda E$ and $q\Lambda E$,
with $pE\ne qE$.  Translate its endpoints by $p^{-1}$ and use
Lemma~\ref{pre:quasiconvex-intersection} to obtain
\[
 \Lambda E\cap\Lambda\bigl((p^{-1}q)E(p^{-1}q)^{-1}\bigr)
 =\Lambda\bigl(E\cap(p^{-1}q)E(p^{-1}q)^{-1}\bigr).
\]
The intersection subgroup is nontrivial cyclic.  If
$p^{-1}q=e_1xe_2$, with $e_1,e_2\in E$ and $x$ a chosen
double-coset representative, it is conjugate by $e_1$ to
$E\cap xEx^{-1}$.  Its axis root is therefore one of the included
classes.  This proves~(ii).

Finally, suppose an axis root $u$ has image $r^k$, where $r$ is a
root in $P$ and $|k|>1$.  Write $E$ for the image subgroup, which in a
noncyclic transition is $I$ or $J$.  We have $r\notin E$, since
otherwise the inverse image of $r$ would exhibit $u$ as a proper
power.  The distinct cosets $E,rE$ both carry the axis of $r$.
Assertion~(ii) proves~(iii).
\end{proof}

We next record exactly how a repetition in this finite set produces a
group-theoretic obstruction.  The observation applies even when the
traversal signs change between the two repeated positions.

\begin{lemma}\label{stat:repetition}
Let $(h_s e)$ be a reduced edge path in $\mathcal T$, indexed by
consecutive integers, and let $r_s\in P$ be nontrivial.  Suppose the
cyclic subgroups
\[
                 h_s\langle r_s\rangle h_s^{-1}
\]
at consecutive positions are commensurable.  If $i<j$ have the same
traversal sign and $\langle r_i\rangle,\langle r_j\rangle$ are
conjugate in $P$, then $G$ contains a Baumslag--Solitar subgroup.
\end{lemma}
\begin{proof}
Choose $p\in P$ with
$p\langle r_i\rangle p^{-1}=\langle r_j\rangle$, and put
$g=h_jph_i^{-1}$.  Since $p$ fixes $e$, this element takes the edge
$h_i e$ to $h_j e$.  It takes the cyclic subgroup at position $i$ to
the subgroup at position $j$, which is commensurable with the first
one.

Let $\sigma$ be the segment from the midpoint of $h_i e$ to the
midpoint of $h_j e$, in the order of the given path.  At each join of
two consecutive translates $g^k\sigma$, the incoming and outgoing
segments use opposite halves of the same edge.  Indeed, $g$ preserves
the canonical edge orientation, and the two positions have the same
traversal sign.  Thus these translates concatenate without
backtracking to a bi-infinite geodesic.  The element $g$ translates
that geodesic a positive distance and is tree-hyperbolic.

The cyclic subgroup at position $i$ is infinite and elliptic.
Lemma~\ref{pre:elliptic-commensurator} applies to it and to $g$, and
gives the asserted embedded Baumslag--Solitar subgroup.
\end{proof}

\begin{proposition}\label{stat:no-atoms}
Let $(\mu_s)_{s\in\mathbb Z}$ be a bounded sequence of nonzero
$P$-currents with transition currents and compatible essential types
as in Proposition~\ref{ann:trajectory}.  If $G$ is Baumslag--Solitar
free, every $\mu_s$ is atomless.
\end{proposition}
\begin{proof}
Suppose a row has an atom, and start indexing there.  At each forward
transition, the coset formula for the induced current has a positive
contribution to this atom from an atom of the transition current.
Choose one such atom and take its image in the next row.  Local
finiteness of the coset sum and injectivity of the boundary maps in
Lemma~\ref{ann:current-induction} justify this choice.
Lemma~\ref{stat:atomic-axis} makes every chosen
line a periodic axis.  Let $r_s$ be its actual root in the coordinates
of the $s$th edge.

We first develop these matches into an actual reduced tree path,
so that the repetition criterion of Lemma~\ref{stat:repetition}
will apply.  For a transition
with representative $d$, use the arriving and departing embeddings
$f,l:P\to H$ from the annulus description.  A generator $k\in K_\tau$
of the stabilizer of the chosen upstairs axis has images of the form
\[
 \theta_{\tau,0}(k)=p^{-1}r_s^{a}p,\qquad
 \theta_{\tau,1}(k)=q^{-1}r_{s+1}^{b}q,
 \qquad a,b\in\mathbb Z\setminus\{0\}.
\]
The bridge
\[
                         f(p)d\,l(q)^{-1}
\]
therefore matches $r_s^a$ with $r_{s+1}^b$.  It belongs to the same
double coset as $d$, so it preserves essentiality at a reversal.
Develop these bridges using the frame recursion
\eqref{ann:frame-equation}.  They give a reduced path with the prescribed
signs, and the framed cyclic subgroups are commensurable at every
step.  This construction works for arbitrary nonzero exponents,
including negative ones.  It therefore suffices to find two positions
with the same root class and traversal sign.

Call a transition exceptional for this chosen trajectory if its group
is cyclic, if either chosen image line is carried by more than one
coset of its image subgroup, or if an upstairs root has a
proper-power image.  By Lemma~\ref{stat:finite-roots}, at least one
of the adjacent roots at every exceptional transition belongs to
$\mathcal D$.  If infinitely many transitions are exceptional, there
are infinitely many such root occurrences.  Some root class and some
traversal sign occur together at two different positions.

If there are only finitely many exceptional transitions, discard an
initial segment.  On the remaining tail there is exactly one
contributing coset on each side of every push formula.  Hence both
image atoms have exactly the mass of their common upstairs atom.
The chosen atoms have one fixed positive mass on this tail.  The norm
bound and Lemma~\ref{stat:atomic-axis} give a uniform bound on
$|r_s|_{\mathrm{cyc}}$.  There are finitely many conjugacy classes of
roots of bounded cyclic length, so once again a root class repeats at
two positions with the same traversal sign.

In either case, Lemma~\ref{stat:repetition} gives a Baumslag--Solitar
subgroup, a contradiction.
\end{proof}

\subsection{A stationary current}

We now have an atomless bounded trajectory.  The next argument turns
it into a stationary one; the case in which the total mass is summable
must be treated separately from averaging.

\begin{proposition}\label{stat:stationary}
Under the assumptions of Proposition~\ref{stat:no-atoms}, there exist
a nonzero current $\nu$ on $I$ and a nonzero atomless current $\mu$ on
$P$ such that
\begin{equation}\label{stat:stationary-equality}
                         i_*\nu=j_*\nu=\mu.
\end{equation}
In particular, if the restricted HNN extension is not hyperbolic,
Proposition~\ref{ann:trajectory} supplies such a stationary current.
\end{proposition}
\begin{proof}
A nonzero current on an infinite cyclic group induces an atomic
current.  Thus no transition in the atomless trajectory is cyclic.
Every transition is consequently aligned, all traversal signs are
the same, and its maps are $i,j$ or $j,i$ according to that sign.
Reverse the indexing if necessary.  We then have currents $\nu_s$ on
the same group $I$ satisfying
\[
                    i_*\nu_s=\mu_s,
             \qquad j_*\nu_s=\mu_{s+1}.
\]
In particular $I$ is nontrivial.  Fix a free basis of $I$ for its
current norm, and define
\[
 S_N=\sum_{s=-N}^N\nu_s,\qquad
 a_N=\|S_N\|=\sum_{s=-N}^N\|\nu_s\|.
\]
The last equality uses positivity.  Telescoping gives
\begin{equation}\label{stat:telescoping}
             j_*S_N-i_*S_N=\mu_{N+1}-\mu_{-N}.
\end{equation}

First suppose $a_N\to\infty$.  By compactness of the norm-one
current space, pass to a subsequence on which $S_N/a_N$ converges to a
current $\nu$ of norm one.  A uniform norm bound controls the mass of
any fixed compact cylinder: such a cylinder is contained in finitely
many translates of the set of lines through the identity.  The right
side of~\eqref{stat:telescoping}, divided by $a_N$, therefore tends
to zero on every compact cylinder.  Continuity of the induced-current
maps yields $i_*\nu=j_*\nu$.  Their positive lower norm bounds show
that this common current $\mu$ is nonzero.

Now suppose $(a_N)$ is bounded.  The positive bilateral sum
\[
                         \nu=\sum_{s\in\mathbb Z}\nu_s
\]
has finite positive norm and defines a Radon measure: every compact
line set is covered by finitely many translates of the identity line
set.  Monotone convergence permits induction through this sum.
Reindexing the resulting positive series gives
\[
 i_*\nu=\sum_{s\in\mathbb Z}\mu_s
        =\sum_{s\in\mathbb Z}\mu_{s+1}=j_*\nu.
\]
These currents have finite norm by the upper norm bound for
induction, and are nonzero.  Equivalently, summability and norm
comparison make the endpoint norms in~\eqref{stat:telescoping}
tend to zero.

In either case we must still prove atomlessness: a weak limit of
atomless currents can have atoms.  Use the constant bilateral
trajectory with row $\mu$ and transition current $\nu$.  Its
transitions are aligned and positive, it is bounded, and
Proposition~\ref{stat:no-atoms} applies.  Thus $\mu$ is atomless.
Every atom of $\nu$ would contribute positive atomic mass through an
injective boundary map in the nonnegative coset formula.  Thus $\nu$
is atomless as well.
\end{proof}

\subsection{Transfers on a conull set of lines}

Fix the currents in~\eqref{stat:stationary-equality}, and let
$\Omega=\operatorname{supp}\mu\subset\partial^2P$ be the closed
support in line space.  A boundary point is \emph{rational} if a
nontrivial element of the free group fixes it.  Equivalently, its ray
labels are eventually periodic.  We first remove all lines with a
rational endpoint.

\begin{lemma}\label{stat:irrational}
For an atomless free-group current, almost every line has two
irrational endpoints.  An injection between finitely generated free
groups preserves and reflects rationality of boundary points.
\end{lemma}
\begin{proof}
Fix a rational endpoint $\xi$, and orient an axis toward it.  Every
oriented line with positive endpoint $\xi$, other than this axis,
merges with the axis at a first vertex.  Fix a possible merge vertex
$v$, and let $D_v$ be this set of lines.  If $r$ is the axis root
directed toward $\xi$, the sets $r^{-n}D_v$, $n\ge0$, are pairwise
disjoint.  They all pass through $v$, because their merge vertices
are $r^{-n}v$.  They have equal measure by invariance, so finiteness
of the measure on the compact line set through $v$ gives
$\mu(D_v)=0$.  There are countably many merge vertices and rational
endpoints.  Atomlessness removes the exceptional axes, and flip
invariance treats negative endpoints.

An injection preserves rationality by equivariance of its boundary
map.  To prove reflection, identify its image with a finitely
generated subgroup $E$ of a free group.  Suppose $\xi\in\Lambda E$
is fixed by a nontrivial element $r$ of the ambient free group.
Orient $r$ toward $\xi$, and choose an axis vertex $x$.  The vertices
$r^n x$ on the resulting axis ray are uniformly
close to $E$, by quasiconvexity and the fact that $\xi\in\Lambda E$.
Choose $e_n\in E$ at uniformly bounded distance from $r^n x$.
Among the finitely many possibilities for
$(r^n x)^{-1}e_n$, two agree, say at $n>m$.  It follows that
\[
                        e_ne_m^{-1}=r^{n-m}\in E.
\]
The inverse image of this nontrivial element fixes the original
boundary point, proving reflection.
\end{proof}

We write $\partial^2\Lambda E$ for the ordered pairs of distinct
points of $\Lambda E$.  If $E=I$ or $J$, an irrational endpoint
belongs to at most one translate $p\Lambda E$, with translates
indexed by left cosets $pE$.  Indeed, for two distinct cosets the
intersection of the corresponding conjugate subgroups is at most
cyclic, and Lemma~\ref{pre:quasiconvex-intersection} identifies the
intersection of their limit sets.  Every point of that intersection
would be rational.

\begin{proposition}\label{stat:transfers}
There is a translation- and flip-invariant conull Borel set
$\mathcal U\subset\Omega$ with the following properties.
\begin{enumerate}[label=\textup{(\roman*)}]
\item Every line of $\mathcal U$ has two irrational endpoints and is
carried by a unique $I$ coset and a unique $J$ coset.
\item On lines carried by $p\Lambda I$, translate by $p^{-1}$ and
apply the boundary isomorphism from $I$ to $J$.  This forward
transfer, and the analogous backward transfer from $q\Lambda J$,
preserve $\mu$ on their domains.
\item Every defined transfer, its inverse, and every output frame
translation take lines of $\mathcal U$ to lines of $\mathcal U$.
Consequently all successive forward and backward transfers exist on
$\mathcal U$.
\end{enumerate}
Measure preservation in~\textup{(ii)} holds on the full Borel domain
of lines with irrational endpoints, not only after restriction to
$\mathcal U$.
\end{proposition}
\begin{proof}
Let $X\subset\partial^2P$ consist of lines with two irrational
endpoints.  It is a conull Borel set by Lemma~\ref{stat:irrational}.
The coset formulas for $i_*\nu$ and $j_*\nu$ show that almost every
line of $X$ is carried by some $I$ coset and some $J$ coset.
The preceding uniqueness observation proves~(i) wherever the
carriers exist.

Let $\vartheta:\partial I\to\partial J$ denote the boundary
isomorphism induced by $j$, with $I$ identified with its included
image under $i$.  On the Borel domain
\[
                   D_p=X\cap\partial^2(p\Lambda I)
\]
define $F_p(L)=\vartheta(p^{-1}L)$.  Rationality preservation and
reflection give $F_p(D_p)=X\cap\partial^2\Lambda J$, a Borel set.
For every Borel $D\subset D_p$, the
two push formulas each have just one contributing coset on the
indicated sets, and give
\begin{equation}\label{stat:branch-measure}
                    \mu(D)=\nu(p^{-1}D)=\mu(F_pD).
\end{equation}
The middle term is expressed in the boundary coordinates of $I$.
These maps are homeomorphisms between their domains and images in
the corresponding limit-set spaces, so the images in this formula
are Borel.  The same argument applies to the backward branches.
Inverse branches preserve measure as well.  Output translations and
flip preserve it by the definition of a current.  Different
representatives of one coset merely insert an output translation.

For completeness, we now choose a single conull set on which all
these operations can be iterated.  Start in $X$ with the Borel null
set
\[
 \begin{split}
 B={}&(X\setminus\Omega)\\
    &\cup\{L\in X:L\text{ has no }I\text{-coset carrier}\}\\
    &\cup\{L\in X:L\text{ has no }J\text{-coset carrier}\}.
 \end{split}
\]
The complement of the support is null because line space has a
countable base.  Saturate $B$ under the partial forward and backward
branches, their inverses, translations, and flip.  There are
countably many such maps and countably many finite compositions.
Equation~\eqref{stat:branch-measure} shows that each composition
takes null sets to null sets on its domain.  Their saturation is
therefore Borel and null.  Its complement in $X$ is the required
$\mathcal U$.  Including $X\setminus\Omega$ among the initial
discarded sets ensures that all these iterates remain in the
support.
\end{proof}

The set $\mathcal U$ is dense in $\Omega$, but it has not been
asserted to be closed.  Until a separate support argument is given,
statements about successive transfers apply to $\mathcal U$.

\subsection{Actual paths and finite windows}

We next describe precisely which tree path is followed by a line in
$\mathcal U$.  This also connects the current to the finite-rank
window stabilizers that will be used in the next section.

\begin{proposition}\label{stat:windows}
Every line $L\in\mathcal U$ determines an actual aligned bilateral
path of positively traversed edges
\[
                     \cdots,E_{-1},E_0=e,E_1,\cdots
\]
in $\mathcal T$, with transferred lines in its successive edge frames.
The following properties hold.
\begin{enumerate}[label=\textup{(\roman*)}]
\item The path is independent of coset representatives and of
compatible changes of frames and line coordinates.
\item For every finite window containing $E_0$, its stabilizer
$K\le P$ is finitely generated and noncyclic, and the endpoints of
$L$ lie in $\Lambda K$.
\item Two distinct actual aligned windows of the same shape about
$e$ have at most cyclic common stabilizer.  In particular, an
irrational point can belong to the limit sets of at most one such
window stabilizer.
\item If an aligned one-sided window joins $e$ to $he$, its
stabilizer is $S=P\cap hPh^{-1}$.  The boundary map of
\[
                  S\longrightarrow h^{-1}Sh,
                         \qquad x\longmapsto h^{-1}xh,
\]
implements its line transfer.  On the irrational-endpoint domain it
preserves $\mu$; it takes its domain in $\mathcal U$ into
$\mathcal U$.
\end{enumerate}
In~\textup{(iii)}, a shape specifies the numbers of edges before and
after the central edge, all traversed positively.
\end{proposition}
\begin{proof}
In a frame $h$ for the current edge, let $pI$ be the unique forward
carrier coset for its line coordinates.  Choose the next edge to be
$hpt e$, and use the forward transfer for its line coordinates.  For
the backward extension, the unique coset $pJ$ gives $hpt^{-1}e$.
Proposition~\ref{stat:transfers} permits indefinite iteration in
both directions.  The frame description in Lemma~\ref{ann:frames}
shows that these are aligned positive paths when ordered from left
to right.  In particular they are reduced.

Changing $p$ to $pi_0$ in a forward step, with $i_0\in I$, does not
change the next edge, since
\[
                       pi_0t=pt\,j(i_0).
\]
It changes its frame on the right by $j(i_0)\in P$ and its line
coordinates by the inverse translation.  Subsequent carrier cosets
change by the same coordinate translation, leaving subsequent
actual edges unchanged.  The backward calculation uses $J$ and the
inverse isomorphism.  Arbitrary changes of an output frame by $P$
have the same effect.  For example, after a forward step to frame
$hpt$, the new line lies in $\partial^2\Lambda J$, so its backward
choice is the identity $J$ coset.  The backward edge is then
$hptt^{-1}e=hpe=he$, with the original line translated into frame
$hp$.  This proves~(i), including compatibility of the forward and
backward constructions.

To prove~(ii), build a finite window one end at a time.  Initially
its stabilizer is $P$.  In the frame at the end being extended, let
the current stabilizer be $K'$.  A forward extension intersects it
with $pIp^{-1}$, and a backward extension intersects it with
$pJp^{-1}$.  Both subgroups are finitely generated subgroups of a
free group.  Their intersection is finitely generated, and
Lemma~\ref{pre:quasiconvex-intersection} says that its limit set is
the intersection of their limit sets.  The current endpoints belong
to both, by induction and by the carrier choice.  They therefore
belong to the new stabilizer limit set.  Conjugation by the change of
frame transfers this assertion to the next coordinates; its
boundary map restricts compatibly to every smaller subgroup.
The same intersection computation, without the endpoint assertion,
proves finite generation for every finite aligned window.  Induction
proves the endpoint assertion in the central coordinates.  A
trivial or cyclic subgroup cannot carry an irrational endpoint, so
the stabilizer is noncyclic.

For~(iii), Lemma~\ref{ann:distinct-paths} gives the cyclic common
stabilizer at the first divergence in either direction.  The window
stabilizers are finitely generated, so
their limit sets intersect in the limit set of their intersection.
This proves the assertion about irrational endpoints as well.
The same proof compares restrictions of longer paths to any common
finite shape.

An element fixing the two endpoint edges of a tree segment fixes
every edge between them.  Thus the stabilizer in~(iv) is precisely
$P\cap hPh^{-1}$.  Successive coordinate changes along the path
compose to conjugation by $h^{-1}$.  Their boundary maps agree with
the boundary map on this subgroup, giving the asserted formula.
If an irrational line has its endpoints in $\Lambda S$, every
successive carrier needed by this prescribed path exists.  Its
uniqueness forces the successive branch choices to be exactly
those of the path.  The measure-preserving branch identities in
Proposition~\ref{stat:transfers} therefore compose to measure
preservation for this partial conjugation.  The invariance of
$\mathcal U$ gives the last assertion.
\end{proof}

\begin{lemma}\label{stat:window-finiteness}
Fix a finite positive window shape about $e$ and a finite ball in the
Cayley tree of $P$.  There are only finitely many actual windows of
that shape whose stabilizer is noncyclic and whose stabilizer limit
set has convex hull meeting that ball.
\end{lemma}
\begin{proof}
The window bound, Lemma~\ref{ann:window-bound}, gives finitely many
$P$-orbits of such windows.  Choose representatives $W_1,\ldots,W_k$
with stabilizers $K_1,\ldots,K_k$.  For each $i$, the quotient of
$\operatorname{Hull}(\Lambda K_i)$ by $K_i$ is a finite core
graph, since $K_i$ is finitely generated.  Choose finitely many
vertices $v_{i,j}$ representing its vertex orbits.

If the hull associated to $pW_i$ meets the fixed ball, enlarge that
ball by one edge if necessary and choose a vertex $x$ in the
intersection.  Then $p^{-1}x=kv_{i,j}$ for some $k\in K_i$ and
some $j$.  Consequently
\[
                     pW_i=xv_{i,j}^{-1}W_i,
\]
because $k$ stabilizes $W_i$.  There are finitely many choices of
$x,i,j$, proving the claim.
\end{proof}

We have obtained a nonzero stationary current and a conull family of
unique bilateral trajectories.  Every finite part of such a
trajectory has a noncyclic stabilizer of the uniformly bounded
total reduced rank from Lemma~\ref{ann:window-bound}.  The next
section encodes these stabilizers by finite core graphs.  Passing
to the closed support there will use closedness of finite-graph
liftability, rather than treating a conull assertion as a support
assertion.
